% ============================================================================
% III. GIỚI THIỆU BÀI TOÁN VÀ CÁC PHƯƠNG PHÁP LIÊN QUAN
% ============================================================================
\section{GIỚI THIỆU BÀI TOÁN VÀ CÁC PHƯƠNG PHÁP LIÊN QUAN}

\subsection{Bối cảnh và động lực}
Các hệ thống Retrieval-Augmented Generation [1] truyền thống thường vận hành theo một đường ống tĩnh một chiều và bộc lộ hai điểm nghẽn căn bản. Thứ nhất, ở pha ingestion, việc dùng chung một bộ \texttt{fixed-window chunker} hoặc \texttt{recursive character splitter} cho mọi loại tài liệu gây ra ba hệ quả trực tiếp: (1) cắt ngang table, code block hoặc điều khoản làm hỏng cấu pháp [10], [11]; (2) làm mất chuỗi tiêu đề cha của các mục con có tên chung chung như \texttt{Options}, \texttt{Parameters}, \texttt{Examples} hay các khoản mục trùng số thứ tự khiến \texttt{embedding model} không phân biệt được mục đó thuộc về chủ đề nào [12]; và (3) làm mất ánh xạ dòng 1:1 đối với \texttt{raw text}. Thứ hai, ở pha runtime, mô hình \texttt{one-shot retrieve-then-read} chỉ nhận một câu truy vấn đơn lẻ từ người dùng, chạy hàm tìm kiếm một lần rồi nhồi toàn bộ \texttt{top-K chunks} vào prompt (\texttt{context stuffing}). Cách làm thụ động này thất bại khi lượt tìm đầu tiên chưa trúng từ khóa của tài liệu, hoặc khi câu trả lời đòi hỏi phải lần theo các tham chiếu chéo giữa nhiều mục và nhiều văn bản khác nhau.

Sự xuất hiện của các \texttt{Coding Agent} cho phép thay đổi tận gốc cả hai điểm nghẽn trên bằng cách chuyển từ RAG tĩnh sang RAG động do \texttt{Agent} chủ động điều khiển. Ở đầu vào, thay vì chấp nhận cắt văn bản mù quáng, người dùng chỉ cần yêu cầu \texttt{Coding Agent} khảo sát quy ước định dạng của kho tài liệu mới và viết một bộ \texttt{Parser} tất định chuyên biệt cho miền đó; bộ \texttt{Parser} này chỉ cần viết và kiểm thử một lần duy nhất rồi chạy bằng thuật toán thông thường trên toàn bộ tài liệu mà không tốn chi phí gọi LLM cho từng trang. Ở đầu ra, thay vì nhồi cố định một danh sách kết quả tìm kiếm vào prompt, hệ thống cung cấp cho \texttt{Agent} một bộ công cụ điều tra để \texttt{Agent} tự lên chiến lược tìm kiếm lặp, tự lật mở cây mục lục, tự đọc nguyên văn các dòng lân cận và tự đan các mối quan hệ vào \texttt{Knowledge Graph} tùy theo năng lực suy luận của mình. Để sự tự chủ đó của \texttt{Agent} ở cả hai đầu luôn bám sát sự thật của tài liệu gốc mà không phát sinh ảo giác, điều cần thiết là một \textbf{RAG framework đứng bên ngoài} định nghĩa ranh giới hợp đồng dữ liệu và cung cấp sẵn hạ tầng lưu trữ, truy xuất lẫn kiểm chứng:
\begin{enumerate}
    \item \textbf{\texttt{Parser interface} với 5 output invariants:} Framework loại bỏ hoàn toàn \texttt{generic chunking}. Mọi bộ \texttt{Parser} do \texttt{Coding Agent} viết muốn lắp vào framework đều phải vượt qua cổng kiểm tra tự động với 5 invariants bắt buộc: đường dẫn phân cấp \texttt{ltree} hợp lệ, khoảng dòng $[s_i, e_i]$ khớp 1:1 với \texttt{raw text}, văn bản nguyên gốc \texttt{verbatim\_text}, văn bản đã ghép chuỗi tiêu đề cha \texttt{contextualized\_text} và mảng \texttt{ancestor\_titles}.
    \item \textbf{Tách biệt biểu diễn lúc indexing và lúc trả kết quả:} Framework tạo \texttt{embedding vector} và \texttt{FTS index} trên trường \texttt{contextualized\_text} để câu truy vấn khớp được cả tên tài liệu/mục cha lẫn nội dung mục con [12]. Ngược lại, khi trả kết quả về cho \texttt{Agent}, framework tách chuỗi tiêu đề cha ra làm thông tin chung của cả \texttt{Provision\allowbreak Group}, chỉ trả \texttt{verbatim\_\allowbreak text} cho từng \texttt{chunk} và sắp xếp các \texttt{chunk} trong nhóm theo \texttt{reading order} tăng dần ($s_i\ \text{ASC}$) kèm nhãn độ tin cậy $\gamma$ từ cổng \texttt{Sparse Match Gating} làm tín hiệu dẫn đường cho bước suy luận tiếp theo của \texttt{Agent}.
    \item \textbf{Bộ 6 MCP tools cho phép \texttt{Agent} chủ động điều tra và kiểm chứng trích dẫn trên PostgreSQL:} Qua giao thức Model Context Protocol [14] trên nền \texttt{Warm Streamable HTTP}, framework trang bị cho \texttt{Agent} 6 công cụ gồm tìm kiếm kênh đôi (\texttt{hybrid\_\allowbreak search}), duyệt cây mục lục 4 hướng (\texttt{hierarchical\_\allowbreak navigate}), đọc lát cắt dòng nguyên văn (\texttt{get\_\allowbreak raw}), duyệt đồ thị đệ quy (\texttt{graph\_\allowbreak traverse}), tự bổ sung cạnh tri thức có kiểm chứng nguyên văn (\texttt{graph\_\allowbreak curate\_\allowbreak edges}) và chốt câu trả lời có đối chiếu tọa độ trích dẫn trực tiếp trên PostgreSQL (\texttt{submit\_\allowbreak grounded\_\allowbreak answer}) [1], [13].
\end{enumerate}

\subsection{Phát biểu bài toán}
Bài toán thiết kế RAG framework phân cấp và \texttt{Parser interface} được phát biểu tổng quát cho mọi miền tài liệu có cấu trúc như sau:

\begin{itemize}
    \item \textbf{Đầu vào của framework:}
    \begin{enumerate}
        \item Một kho tài liệu có cấu trúc phân cấp $\mathcal{D} = \{d_1, d_2, \dots, d_M\}$. Mỗi tài liệu $d_m$ có mã định danh phân cấp $\text{slug}(d_m) \in \mathcal{T}_{\text{ltree}}$ và mảng các dòng raw text $L(d_m) = (\ell_1, \ell_2, \dots, \ell_{|d_m|})$.
        \item Một bộ \texttt{Parser} $f_{\text{parse}}: \mathcal{D} \rightarrow 2^{\mathcal{C}} \times 2^{\mathcal{E}}$ được \texttt{Coding Agent} hoặc lập trình viên viết cho miền tài liệu $\mathcal{D}$ và đăng ký vào \texttt{Parser Registry}.
        \item Runtime query gồm câu hỏi ngữ nghĩa $q$, verbatim pattern hoặc regex $p$, tập lọc tài liệu $\mathcal{S} \subseteq \mathcal{D}$ và \texttt{ltree} path prefix $\pi \in \mathcal{T}_{\text{ltree}}$.
    \end{enumerate}
    \item \textbf{Output contract 5 invariants mà \texttt{Parser} $f_{\text{parse}}$ phải thỏa mãn:}
    Để framework chấp nhận nạp vào database, mỗi \texttt{chunk} $c_i = (\text{path}_i, \text{text}^{\text{verb}}_i, \text{text}^{\text{ctx}}_i, [s_i, e_i], \text{meta}_i) \in \mathcal{C}$ và mỗi graph edge $e \in \mathcal{E}$ do $f_{\text{parse}}$ xuất ra phải thỏa mãn 5 invariants:
    \begin{enumerate}
        \item \textit{Invariant đường dẫn \texttt{ltree}}: $\text{path}_i \in \mathcal{T}_{\text{ltree}}$ là duy nhất trong tài liệu, bắt đầu bằng prefix $\text{slug}(d_m)$ và mỗi nhãn phân cấp khớp regex $\texttt{\^{}[A-Za-z0-9\_]\{1,255\}\$}$.
        \item \textit{Invariant vị trí dòng 1:1}: Chỉ số dòng bắt đầu và kết thúc thỏa mãn $1 \le s_i \le e_i \le |d_m|$, trỏ chính xác vào vị trí dòng trong $L(d_m)$ và không cắt ngang các atomic block như code block hay table.
        \item \textit{Invariant tiền tố ngữ cảnh cha}: Trường văn bản ngữ cảnh có dạng chuẩn $\text{text}^{\text{ctx}}_i = \texttt{"[}\text{context}_i\texttt{]\textbackslash n"} \circ \text{text}^{\text{verb}}_i$, trong đó $\text{context}_i$ là chuỗi tiêu đề từ gốc tài liệu đến mục hiện tại nối bằng ký hiệu \texttt{" > "}.
        \item \textit{Invariant danh sách tiêu đề tổ tiên}: Trường $\text{meta}_i$ chứa mảng \texttt{ancestor\_titles} ứng với từng cấp trên $\text{path}_i$ để framework tự động tái tạo các virtual parent node khi tree navigation.
        \item \textit{Invariant citation grounding trên graph edge}: Mọi cạnh quan hệ $e = (c_u, c_v, r, \tau) \in \mathcal{E}$ phải thuộc tập 9 mã quan hệ chuẩn của hệ thống, thỏa mãn $c_u \neq c_v$ và chứa chuỗi trích dẫn $\tau \neq \emptyset$ xuất hiện nguyên văn trong $\text{text}^{\text{verb}}_u$ sau khi chuẩn hóa khoảng trắng.
    \end{enumerate}
    \item \textbf{Đầu ra của framework khi retrieval và navigation:}
    Từ tập $\mathcal{C}$ và $\mathcal{E}$ đã nạp, với mỗi query $(q, p)$, framework trả về cấu trúc hợp nhất $\mathcal{R}(q, p) = (\mathcal{D}_{\text{hit}}, \mathcal{G}, N_{\text{verb}}, \gamma)$ gồm: danh sách tài liệu deduplicated $\mathcal{D}_{\text{hit}}$; danh sách các context group $\mathcal{G} = (g_1, g_2, \dots, g_K)$ gom các \texttt{chunk} có cùng khóa $(\text{slug}(d_m), \text{context})$ theo thứ tự ưu tiên Venn partitioning ($\mathcal{H}_{\text{sem}} \cap \mathcal{H}_{\text{verb}} \rightarrow \mathcal{H}_{\text{sem}} \setminus \mathcal{H}_{\text{verb}} \rightarrow \mathcal{H}_{\text{verb}} \setminus \mathcal{H}_{\text{sem}}$) nhưng sắp xếp các \texttt{chunk} bên trong mỗi nhóm theo natural reading order $(s_i\ \text{ASC}, \text{path}_i\ \text{ASC})$; tổng số verbatim matches $N_{\text{verb}}$; và confidence label $\gamma \in \{\texttt{none}, \texttt{low}, \texttt{medium}, \texttt{high}\}$, cùng các phép duyệt cây $\mathcal{T}_{\text{ltree}}$, duyệt đồ thị $G = (\mathcal{C}, \mathcal{E})$ và đọc raw line slice.
    \item \textbf{Các câu hỏi nghiên cứu trọng tâm:} Từ phát biểu bài toán trên, thực nghiệm của báo cáo tập trung trả lời ba câu hỏi nghiên cứu:
    \begin{itemize}
        \item \textbf{CH1 (Bảo toàn cấu trúc đa miền qua \texttt{Parser interface} 5 invariants):} Liệu ba bộ \texttt{Parser} chuyên biệt cho Next.js MDX, Node.js Markdown và Văn bản Pháp luật Bảo hiểm tiếng Việt lắp vào framework có phân rã trọn vẹn kho tài liệu mà bảo toàn 100\% ánh xạ dòng 1:1, tiền tố ngữ cảnh phân cấp và ranh giới khối cú pháp hay không?
        \item \textbf{CH2 (Đóng góp của Dual-Channel Venn và giới hạn cửa sổ Reranker):} Sự kết hợp giữa nhánh ngữ nghĩa (Dense + Sparse FTS qua RRF trên $\text{text}^{\text{ctx}}_i$), tầng Cross-Encoder 256 token và kênh nguyên văn regex/trigram đóng góp thế nào vào độ chính xác truy xuất trên hai lớp truy vấn \texttt{TS1-Lexical-API} và \texttt{TS2-Conceptual}?
        \item \textbf{CH3 (Điều tra đa bước qua bộ 6 MCP tools và kiểm chứng trích dẫn tất định):} Khi thay thế lối nhồi ngữ cảnh một chiều bằng bộ 6 MCP tools trên nền \texttt{Warm Streamable HTTP}, \texttt{Coding Agent} phối hợp tìm kiếm kênh đôi, duyệt cây \texttt{ltree}, đọc lát cắt dòng nguyên văn, tự mở rộng \texttt{Knowledge Graph} và kiểm chứng tọa độ trích dẫn trên PostgreSQL như thế nào để giải quyết các truy vấn \texttt{multi-hop} với độ trễ thời gian thực và ngăn chặn trích dẫn giả mạo?
    \end{itemize}
\end{itemize}

\subsection{Các phương pháp liên quan}
Các công trình nghiên cứu nền tảng liên quan đến hybrid retrieval, structural chunking và Knowledge Graph được chia thành bốn nhóm và tổng hợp tại Bảng \ref{tab:related_work}:

\begin{itemize}
    \item \textbf{Lexical retrieval, Dense retrieval và Two-stage hybrid fusion:} Phương pháp lexical matching như BM25 [2] khớp chính xác các ký hiệu hiếm nhưng không xử lý được query khác từ vựng bề mặt. Ngược lại, kiến trúc Bi-Encoder như DPR [3] và embedding model dựa trên Decoder-only Transformer như Qwen3-Embedding [4] kết hợp Matryoshka Representation Learning [5] biểu diễn tương đồng ngữ nghĩa trong không gian liên tục $\mathbb{R}^d$. Công thức Reciprocal Rank Fusion [6] cho phép kết hợp thứ hạng từ hai kênh mà không phụ thuộc thang điểm tuyệt đối, trước khi đưa top candidates qua mô hình Cross-Encoder [7], [8] tính cross-attention đầy đủ giữa từng cặp token.
    \item \textbf{Retrieval tài liệu kỹ thuật hỗ trợ lập trình:} DocPrompting [9] chứng minh việc cung cấp đúng đoạn tài liệu chính thức giúp code generation model gọi chính xác các thư viện mới ngoài tập huấn luyện. Tiếp đó, benchmark CodeRAG-Bench [10] chỉ ra rằng tài liệu kỹ thuật chứa nhiều nhiễu định dạng và các phương pháp single-channel retrieval thường xếp hạng sai khi có nhiều hàm mang chữ ký hoặc tiêu đề mục con tương tự nhau.
    \item \textbf{AST chunking và hierarchical contextual retrieval:} Để tránh cắt đứt code block như fixed-window chunking, cAST [11] chunk mã nguồn theo node boundaries trên Abstract Syntax Tree. Ở bước indexing, phương pháp Contextual Retrieval [12] gắn thêm context prefix của tài liệu vào đầu mỗi \texttt{chunk} trước khi tạo embedding vector và BM25 index, giúp disambiguate các đoạn văn bản phụ thuộc vào tiêu đề mục cha.
    \item \textbf{Knowledge Graph trong RAG và giao thức MCP cho Agent:} GraphRAG [13] xây dựng entity-relation graph để giải quyết các multi-hop queries, trong khi giao thức Model Context Protocol [14] chuẩn hóa interface để các \texttt{Agent} gọi công cụ retrieval và thao tác dữ liệu có cấu trúc.
\end{itemize}

\begin{table}[H]
\centering
\caption{Tổng hợp các công trình nghiên cứu liên quan và giới hạn khi áp dụng lên tài liệu phân cấp phức tạp}
\label{tab:related_work}
\footnotesize\setlength{\tabcolsep}{4pt}\begin{tabularx}{\textwidth}{>{\raggedright\arraybackslash}p{2.9cm}>{\raggedright\arraybackslash}p{3.1cm}>{\raggedright\arraybackslash}p{3.4cm}>{\raggedright\arraybackslash}X}
\toprule
\textbf{Công trình} & \textbf{Kỹ thuật cốt lõi} & \textbf{Phạm vi giải quyết chính} & \textbf{Giới hạn trên tài liệu phân cấp và tài liệu kỹ thuật} \\
\midrule
Robertson \& Zaragoza (2009) -- BM25 [2] & TF-IDF term weighting với length saturation & Khớp chính xác từ khóa hiếm, tên định danh và mã lỗi & Không xử lý được paraphrased query khác từ vựng bề mặt; standard stemmer làm biến dạng định danh chuyên ngành \\
\addlinespace[3pt]
Karpukhin et al. (2020) [3]; Zhang et al. (2025) [4]; Kusupati et al. (2022) [5] & Bi-Encoder Dense Embedding và Matryoshka dimension truncation & Biểu diễn ngữ nghĩa câu hỏi tự nhiên và văn bản dài trong không gian $\mathbb{R}^d$ & Dễ làm mờ khác biệt nhỏ giữa các định danh gần giống nhau; luôn trả về nearest neighbors ngay cả với out-of-domain queries \\
\addlinespace[3pt]
Cormack et al. (2009) [6]; Nogueira \& Cho (2019) [7]; Wang et al. (2021) [8] & RRF rank fusion kết hợp Cross-Encoder reranking & Cân bằng độ phủ lexical -- semantic và tinh chỉnh thứ tự top-$K$ bằng cross-attention & Context window của Cross-Encoder nhỏ gọn thường ngắn hơn các mục tài liệu lớn; thiếu kênh regex matching \\
\addlinespace[3pt]
Zhou et al. (2023) -- DocPrompting [9]; Wang et al. (2025) -- CodeRAG-Bench [10] & Retrieval tài liệu API chính thức đưa vào context sinh code & Chứng minh hiệu quả của tài liệu hàm đối với các API mới hoặc ngoài tập huấn luyện & Coi tài liệu là các đoạn văn phẳng có sẵn; chỉ thực hiện \texttt{one-shot retrieval} mà không cung cấp công cụ duyệt cấu trúc cho \texttt{Agent} \\
\addlinespace[3pt]
Zhang et al. (2025) -- cAST [11]; Anthropic (2024) -- Contextual Retrieval [12] & AST chunking và context prefixing vào đầu chunk khi indexing & Bảo toàn ranh giới khối cú pháp và disambiguate các đoạn con trùng tiêu đề & Hardcode \texttt{Parser} cho mã nguồn thuần túy thay vì tài liệu lai; context prefix bị lặp lại dư thừa khi trả nhiều chunk cùng mục \\
\addlinespace[3pt]
Edge et al. (2024) -- GraphRAG [13] & Xây dựng entity-relation graph phục vụ multi-hop queries & Liên kết thông tin phân tán giữa nhiều tài liệu khác nhau & Graph sinh tự động bằng LLM ở pha nạp gây tốn kém và dễ phát sinh hallucinated edges thiếu citation grounding \\
\bottomrule
\end{tabularx}
\end{table}

\noindent\textbf{Technical Gap:} Từ Bảng \ref{tab:related_work}, có thể thấy các giải pháp hiện có vẫn bị bó hẹp trong mô hình RAG tĩnh: ở pha nạp thì dùng \texttt{generic chunker} làm hỏng cấu trúc hoặc gọi LLM sinh ngữ cảnh và đồ thị cho từng trang gây tốn kém [12], [13]; ở pha truy vấn thì nhồi thụ động kết quả của một lượt tìm kiếm duy nhất vào prompt nên dễ thất bại trước các câu hỏi \texttt{multi-hop}. Khoảng trống đặt ra là xây dựng một \textbf{RAG framework động định hướng \texttt{Coding Agent} ở cả hai đầu}: (1) cung cấp \texttt{Parser interface} 5 invariants để \texttt{Coding Agent} tự viết \texttt{Parser} tất định cho từng miền dữ liệu mới, và (2) cung cấp bộ 6 MCP tools trên tầng PostgreSQL đa index để \texttt{Coding Agent} chủ động điều hướng cây \texttt{ltree}, đọc \texttt{raw line slice}, tự curate \texttt{Knowledge Graph} và chịu sự kiểm chứng trích dẫn tất định trước khi trả lời [14].

\subsection{Đóng góp của báo cáo}
Trên cơ sở kế thừa mô hình RRF + Cross-Encoder [6]–[8], kỹ thuật Matryoshka dimension truncation [5] và nguyên lý context prefixing [12], báo cáo thực hiện hai nhóm đóng góp chính:
\begin{enumerate}
    \item \textbf{Chuyển dịch kiến trúc từ RAG tĩnh một chiều sang RAG động định hướng \texttt{Coding Agent}:}
    \begin{itemize}
        \item Ở pha ingestion, thiết kế \texttt{Parser interface} với 5 output invariants, bộ hàm chuẩn hóa đường dẫn \texttt{ltree} dùng chung và ingestion pipeline có \texttt{differential vector cache}, cho phép \texttt{Coding Agent} tự viết và lắp ghép các bộ \texttt{Parser} tất định cho miền tài liệu mới mà không cần sửa đổi lõi hệ thống.
        \item Ở lõi truy xuất, xây dựng \texttt{dual-channel retrieval engine} kết hợp Hybrid Search trên vector 512 chiều của \texttt{Qwen3-Embedding-0.6B} [4] và FTS cấu hình \texttt{simple} có proximity boost qua RRF [6], reranking bằng \texttt{mmarco-mMiniLMv2-L12-H384-v1} [8], hợp nhất theo \texttt{Venn partitioning} với kênh tìm kiếm regex/trigram, tách tiêu đề cha ra cấp nhóm theo đúng \texttt{reading order} và xuất nhãn độ tin cậy $\gamma$ qua cổng \texttt{Sparse Match Gating} làm tín hiệu dẫn đường cho \texttt{Agent}.
        \item Ở pha runtime, thay thế lối \texttt{context stuffing} thụ động bằng bộ 6 MCP tools [14] trên nền \texttt{Warm Streamable HTTP} cùng giao diện Web Observatory Dashboard, cho phép \texttt{Coding Agent} chủ động tìm kiếm nhiều bước, duyệt cây \texttt{ltree}, trích xuất dòng nguyên văn, tự curate cạnh vào \texttt{Knowledge Graph} dưới ràng buộc \texttt{citation\_text} và kiểm chứng tọa độ trích dẫn câu trả lời trực tiếp trên PostgreSQL.
    \end{itemize}
    \item \textbf{Hiện thực hóa ba bộ \texttt{Parser} trên Next.js MDX, Node.js Markdown và Văn bản Pháp luật Bảo hiểm Việt Nam để kiểm định framework:}
    \begin{itemize}
        \item Xây dựng ba bộ \texttt{Parser} chuyên biệt giải quyết các cấu trúc phức tạp của tài liệu kỹ thuật tiếng Anh (Next.js, Node.js) và văn bản quy phạm pháp luật tiếng Việt (Luật Kinh doanh bảo hiểm 08/2022/QH15, Nghị định 46/2023/NĐ-CP, Thông tư 67/2023/TT-BTC), indexing 170 tài liệu gồm 171.183 dòng raw text thành 5.975 \texttt{chunk} bảo toàn 100\% code block, điều khoản pháp lý và vị trí dòng 1:1; thiết lập bộ dữ liệu đánh giá 100 truy vấn qua quy trình 3 pha cách ly mù cùng nghiên cứu Ablation, kiểm định thống kê McNemar và kiểm chứng các chu trình điều tra \texttt{multi-hop} bằng bộ 6 MCP tools.
    \end{itemize}
\end{enumerate}
